% 04_data_storage.tex
\section{Data Storage}

\subsection{Data Storage and Input Management}
The system employs a decoupled approach to data management, distinguishing between transient input buffers and persistent structural representations. This logic is primarily orchestrated by the \texttt{DataManager} class.

\subsubsection{\texttt{DataManager} (The Ingestion Hub)}
The \texttt{DataManager} serves as a centralized utility for data ingestion and preprocessing. It does not maintain the long-term state of a data structure but rather prepares the "raw" data required to initialize or modify one.

\begin{itemize}
    \item \textbf{Transient Storage}: Internally, the class utilizes several \texttt{std::vector} containers to hold different types of raw inputs: \texttt{data} (for linear structures and trees), \texttt{dataGraph} (a collection of \texttt{Edge} objects), and \texttt{dataGridGraph} (a 2D matrix for pathfinding).
    \item \textbf{Input Modalities}: The implementation supports three primary input streams:
        \begin{itemize}
            \item \textbf{File I/O}: Methods like \texttt{inputFromFile()} parse external text files, ensuring data persistence across sessions.
            \item \textbf{Random Generation}: A suite of randomizers
            
            (e.g., \texttt{randomDataDAG}, \texttt{randomDataPlanarGraph}) generates synthetically valid test cases with configurable constraints (min/max values, density).
            \item \textbf{Manual Console Input}: Direct user input is captured via string parsing in 
            
            \texttt{inputFromConsole()}.
        \end{itemize}
    \item \textbf{Reset Policy}: To ensure consistency, the \texttt{DataManager} follows a "clear-on-action" policy. Before any new data is ingested, existing buffers are cleared (e.g., \texttt{data.clear()}), preventing data leakage between unrelated visualization steps.
\end{itemize}

\subsubsection{Structure-Specific Persistence}
Once the \texttt{DataManager} has prepared the raw input, the data is passed to a specific \texttt{IVisualizable} object (such as \texttt{AVLTree} or \texttt{MaxHeap}). 

\begin{itemize}
    \item \textbf{Data Ownership}: Upon initialization, the concrete data structure class takes ownership of the data and converts it into its own internal format (e.g., a pointer-based tree of \texttt{Node} objects or a heap-ordered \texttt{std::vector}).
    \item \textbf{State Preservation}: While the \texttt{DataManager} may reset for the next user interaction, the Model layer retains the current state of the structure to facilitate frame-by-frame playback and incremental updates (Insertion/Deletion) until the user explicitly switches the active structure via the \texttt{AppEngine}.
\end{itemize}
